% =====================================================================
%  The Purpose of a Character:
%  Identity as a Conserved Invariant and Purpose as a Fixed Point,
%  for a Single Agent and for Agents in Society
%
%  Self-contained. States its own axioms; re-derives every primitive
%  it uses; cites only standard, externally published mathematics.
% =====================================================================
\documentclass[11pt,twocolumn,a4paper]{article}

\usepackage[utf8]{inputenc}
\usepackage[T1]{fontenc}
\usepackage{lmodern}
\usepackage[margin=1in]{geometry}
\usepackage{amsmath,amssymb,amsthm,mathtools}
\usepackage{thmtools}
\usepackage{enumitem}
\usepackage{microtype}
\usepackage{booktabs}
\usepackage{graphicx}
\usepackage{caption}
\usepackage{hyperref}
\usepackage{cleveref}
\usepackage{xcolor}
\usepackage{tikz}
\usetikzlibrary{arrows.meta,positioning,calc}

\hypersetup{
  colorlinks=true, linkcolor=blue!50!black,
  citecolor=green!50!black, urlcolor=blue!60!black
}

% ---- theorem environments ----
\declaretheoremstyle[
  spaceabove=6pt, spacebelow=6pt,
  headfont=\bfseries, notefont=\mdseries, notebraces={(}{)},
  bodyfont=\itshape, postheadspace=1em
]{thmstyle}
\declaretheoremstyle[
  spaceabove=6pt, spacebelow=6pt,
  headfont=\bfseries, notefont=\mdseries, notebraces={(}{)},
  bodyfont=\normalfont, postheadspace=1em
]{defstyle}

\declaretheorem[style=thmstyle, name=Theorem, numberwithin=section]{theorem}
\declaretheorem[style=thmstyle, name=Proposition, sibling=theorem]{proposition}
\declaretheorem[style=thmstyle, name=Lemma, sibling=theorem]{lemma}
\declaretheorem[style=thmstyle, name=Corollary, sibling=theorem]{corollary}
\declaretheorem[style=defstyle, name=Definition, sibling=theorem]{definition}
\declaretheorem[style=defstyle, name=Axiom, sibling=theorem]{axiom}
\declaretheorem[style=defstyle, name=Remark, sibling=theorem]{remark}
\declaretheorem[style=defstyle, name=Example, sibling=theorem]{example}

\crefname{axiom}{Axiom}{Axioms}
\Crefname{axiom}{Axiom}{Axioms}

% ---- notation ----
\newcommand{\R}{\mathbb{R}}
\newcommand{\N}{\mathbb{N}}
\newcommand{\Agent}{\mathsf{A}}
\newcommand{\Self}{\Gamma}          % the self-graph
\newcommand{\V}{V}
\newcommand{\E}{E}
\newcommand{\wt}{w}
\newcommand{\floorc}{\beta}          % individuation floor
\newcommand{\cut}{\partial}
\newcommand{\Char}{\chi}             % character invariant
\newcommand{\drive}{\gamma}          % drive / gating field
\newcommand{\state}{s}
\newcommand{\Beh}{\mathcal{B}}       % behaviour space
\newcommand{\Purp}{\mathcal{P}}      % purpose cell
\newcommand{\Soc}{\Sigma}            % society (quotient)
\newcommand{\abs}[1]{\lvert#1\rvert}
\newcommand{\norm}[1]{\lVert#1\rVert}
\newcommand{\dist}{d}

\title{\textbf{The Purpose of a Character:\\
Identity as a Conserved Invariant and Purpose as a Fixed Point,\\
for a Single Agent and for Agents in Society}}

\author{Kundai Farai Sachikonye\\[2pt]
\small AIMe Registry for Artificial Intelligence\\
\small \texttt{kundai.sachikonye@bitspark.com}}
\date{\today}

\begin{document}
\maketitle

\begin{abstract}
\noindent
We ask what it is for a synthetic character---an agent in a simulated
world---to \emph{have} an identity and a \emph{purpose}, and we answer
the question with mathematics that a reader can follow without reference
to any other work. We model an agent as a finite weighted graph, the
\emph{self-graph}, whose vertices are the distinctions the agent draws
within itself and whose edge weights are the costs of holding those
distinctions apart. From three axioms---the agent is finite, it is never
in a null state, and no self-distinction is free---we prove a positive
\emph{floor} on every internal separation and derive two objects that
together constitute the agent's character. \textbf{(T1)} Each agent
carries a \emph{conserved invariant}, computed from its own internal
boundary content, that is preserved under every relabelling of its parts
and is realised by no single part; this invariant is the agent's
\emph{identity}, the standing structure from which its changing behaviour
derives, and it is a region, never a point. \textbf{(T2)} An agent's
moment-to-moment conduct is a walk on its self-graph, undetermined by the
graph alone; selecting the next act requires an additional datum, a
\emph{drive field}, and the field's stable set---its fixed point---is the
agent's \emph{purpose}: the region of behaviour the agent is drawn to and
returns toward. \textbf{(T3)} Purpose exists and is unique up to an
equivalence we make precise; \textbf{(T4)} it is an attractor, so conduct
perturbed away from purpose returns to it, giving characters that are
robust rather than brittle. We then place the agent in \emph{society}.
\textbf{(T5)} A collection of agents acting as one admits a single shared
drive on the quotient of their self-graphs, and exactly one; two drives
are the signature of a group that is really two. \textbf{(T6)} Coordinated
purpose among agents requires no shared identity, shared beliefs, or
shared internal vocabulary: agents with wholly disjoint self-graphs can
share a purpose defined in the common space of outcomes, and \textbf{(T7)}
a crowd of such agents settles on a purpose more sharply than any one of
them. Finally \textbf{(T8)} an agent's identity is unreachable from
outside: no finite interrogation resolves it to a point, and an authored
copy of a character is always a distinct individual with its own history.
We give constructions, an attractor analysis, and a numerical study on
explicitly built agents. The interpretive readings---identity as
character, drive as motivation, purpose as goal, the quotient as a
collective---are marked as such and no theorem depends on them.
\end{abstract}

\vspace{0.5em}
\noindent\textbf{Keywords:} agent identity; purpose; conserved invariant;
gated walk; fixed point; attractor; quotient graph; coordination;
non-player character.

% =====================================================================
\section{Introduction}
\label{sec:intro}
% =====================================================================

\subsection{The question}

A character in a simulated world is convincing when it seems to have a
life of its own: a standing identity that persists across encounters, and
a purpose it pursues whether or not anyone is watching. The difficulty is
not to make an agent talk or react---that is comparatively easy---but to
make it \emph{present}: to give it an interior it is not always turning
outward. A merely reactive agent has no identity of its own; it is a
mirror of whoever addresses it. What distinguishes a character from a
reaction is that the character has something conserved inside it that its
behaviour derives from, and something it is trying to do that is its own.

This paper asks what those two things \emph{are}, formally. We want a
definition of a character's \emph{identity} and of its \emph{purpose}
that is precise enough to compute with, general enough to apply to any
bounded agent, and self-contained: the reader should need nothing but
this paper and standard mathematics.

\subsection{Method}

We prove a conditional. \emph{If} an agent is a finite thing that is
always in some state and that cannot draw an internal distinction for
free, \emph{then} it necessarily has an identity of a determinate kind
and admits a purpose of a determinate kind. We assume nothing about what
the agent is made of. The deductions are combinatorial (finite graphs and
their symmetries), order-theoretic (walks and their histories), and
elementary in dynamics (fixed points and attractors of a gradient flow).
Where a tempting claim could not be proved at this standard we have
removed it.

Our strategy throughout is to model the agent as a finite weighted graph
and to read identity, behaviour, and purpose off its structure:

\begin{itemize}[leftmargin=*]
\item the \emph{self-graph} is the agent's internal structure---the
distinctions it holds within itself and the cost of holding them;
\item \emph{identity} is the invariant of that graph that survives all
re-description;
\item \emph{behaviour} is a walk on the graph;
\item \emph{purpose} is the fixed point of the field that selects the
walk.
\end{itemize}

\subsection{What we do and do not claim}

We do not claim that any particular agent architecture is correct, nor
that biological or human identity is a finite graph. We claim that
\emph{any} bounded agent satisfying the three axioms has the structure we
derive, and that this structure is enough to found a usable notion of
character and purpose for synthetic agents. The reading of the formal
objects as ``character,'' ``motivation,'' and ``goal'' is offered as
orientation and is confined to remarks; the theorems are graph- and
order-theoretic and are used only as such.

\subsection{Organisation}

\Cref{sec:setting} fixes the agent as a self-graph and states the axioms.
\Cref{sec:floor} proves the floor. \Cref{sec:identity} derives identity
(T1). \Cref{sec:behaviour} makes behaviour a gated walk and proves that
selection needs a drive (T2). \Cref{sec:purpose} defines purpose as a
fixed point and proves existence, uniqueness, and attraction (T3, T4).
\Cref{sec:society} places agents in society and proves the single-drive,
disjoint-coordination, and crowd-sharpening results (T5--T7).
\Cref{sec:unreachable} proves identity is unreachable from outside and
copies are distinct (T8). \Cref{sec:validation} reports a numerical
study. \Cref{sec:discussion} discusses scope.

% =====================================================================
\section{The agent as a self-graph}
\label{sec:setting}
% =====================================================================

We build the agent from the inside. An agent is not a point; it is a
system of internal distinctions---this mood as against that, this
commitment as against its abandonment, this memory as against its
absence. We record those distinctions as a graph.

\begin{definition}[Self-graph]
\label{def:selfgraph}
The \emph{self-graph} of an agent $\Agent$ is a finite weighted graph
$\Self=(\V,\E,\wt)$ where $\V$ is a finite set of \emph{parts} (the
internal distinctions the agent draws within itself), $\E\subseteq
\binom{\V}{2}$ is a set of \emph{separations} (pairs of parts the agent
holds apart), and $\wt:\E\to(0,\infty)$ assigns to each separation the
\emph{cost} of holding it. The graph is simple and, by
\Cref{ax:present} below, connected.
\end{definition}

\begin{definition}[Cut and boundary]
\label{def:cut}
For disjoint $S,T\subseteq\V$, the \emph{cut} between them is
$\cut(S,T)=\{uv\in\E: u\in S, v\in T\}$ with weight
$\wt(\cut(S,T))=\sum_{e\in\cut(S,T)}\wt(e)$. For $\varnothing\neq
U\subsetneq\V$ the \emph{boundary} of $U$ is $\cut(U)=\cut(U,\V\setminus
U)$ and its \emph{boundary cost} is $b(U)=\wt(\cut(U))$.
\end{definition}

We now state the three axioms. They are meant to be almost
undeniable for anything we would call a bounded agent.

\begin{axiom}[Finiteness]
\label{ax:finite}
The self-graph is finite: $\abs{\V}<\infty$. An agent holds finitely many
internal distinctions at a time.
\end{axiom}

\begin{axiom}[Presence, or no null state]
\label{ax:present}
At every instant the agent occupies exactly one part; there is no
null part and no gap between parts. Equivalently, the self-graph is
connected and nonempty: the agent is always somewhere in itself.
\end{axiom}

\begin{axiom}[Costly individuation]
\label{ax:costly}
No internal distinction is free: every separation has positive cost, and
these costs are bounded below by a constant. Formally, there is
$\floorc>0$ with $\wt(e)\ge\floorc$ for all $e\in\E$.
\end{axiom}

\begin{remark}[Why costly individuation is not an extra assumption]
\label{rem:costly-earned}
\Cref{ax:costly} can be motivated rather than merely posited. To
individuate a part $U$ of the agent is to tell it apart from the rest of
the agent, $\V\setminus U$. If the agent is a genuine whole---if
distinguishing a part against the rest were a task that could be
\emph{finished}, the rest would be exhausted and the agent would be
merely that finite catalogue with nothing held in reserve. An agent that
is always in some state (\Cref{ax:present}) and never reduces to a closed
catalogue leaves, on each of its parts, an irreducible remainder of the
comparison against the rest; that remainder is the positive floor. We do
not rely on this motivation---we take \Cref{ax:costly} as an axiom and
prove everything from it---but it explains why the floor is the natural,
not an arbitrary, hypothesis. The floor is the trace the agent-as-a-whole
leaves on each of its parts.
\end{remark}

\begin{example}[A three-mood agent]
\label{ex:threemood}
Let $\V=\{a,b,c\}$ with separations $ab,bc,ca$ of costs $2,3,2$. This is a
weighted triangle with $\floorc=2$. Every part has boundary cost: e.g.
$b(\{a\})=\wt(ab)+\wt(ca)=4$. We return to this agent throughout.
\end{example}

% =====================================================================
\section{The floor: no internal distinction is free}
\label{sec:floor}
% =====================================================================

The first consequence is immediate but load-bearing: every part of the
agent is separated from the rest at strictly positive cost. Nothing in
the agent is marked off for nothing.

\begin{theorem}[Floor]
\label{thm:floor}
Under \Cref{ax:finite,ax:present,ax:costly}, every nonempty proper part
$U\subsetneq\V$ has boundary cost bounded below by the floor:
\[
b(U)=\wt(\cut(U))\ \ge\ \floorc\ >\ 0.
\]
No part is separated from the rest of the agent at zero cost.
\end{theorem}

\begin{proof}
By \Cref{ax:present} the self-graph is connected, so for a nonempty
proper $U$ there is at least one edge from $U$ to $\V\setminus U$;
otherwise $U$ and its complement would be disconnected. Hence
$\cut(U)\neq\varnothing$, and $b(U)=\sum_{e\in\cut(U)}\wt(e)\ge\wt(e^\ast)
\ge\floorc>0$ for any $e^\ast\in\cut(U)$ by \Cref{ax:costly}.
\end{proof}

\begin{corollary}[The realised floor is attained]
\label{cor:realised}
The quantity $\floorc_\Self:=\min_{\varnothing\neq U\subsetneq\V}b(U)$ is
strictly positive and attained, by finiteness of $\V$.
\end{corollary}

\begin{proof}
Positivity is \Cref{thm:floor}; attainment is minimisation of a
real-valued function over the finite nonempty set of proper parts.
\end{proof}

% =====================================================================
\section{Identity: the conserved invariant}
\label{sec:identity}
% =====================================================================

We now isolate what, in an agent, stays fixed while everything on its
surface changes. The surface is the labelling of parts---which mood is
called which, in what order the agent's states are enumerated. A
re-description of the agent that preserves which parts are separated from
which, and at what cost, changes nothing real about the agent; it only
relabels. We define this precisely and find the quantity that survives
all such relabelling.

\subsection{Re-description is relabelling}

\begin{definition}[Isomorphism; automorphism]
\label{def:iso}
A \emph{weighted isomorphism} $\phi:\Self\to\Self'$ is a bijection
$\phi:\V\to\V'$ with $uv\in\E\iff\phi(u)\phi(v)\in\E'$ and
$\wt'(\phi(u)\phi(v))=\wt(uv)$. An \emph{automorphism} is a weighted
isomorphism $\Self\to\Self$; the automorphisms form a group
$\mathrm{Aut}(\Self)$. A quantity $\iota$ is an \emph{invariant} if
$\iota(\Self)=\iota(\Self')$ whenever $\Self\cong\Self'$.
\end{definition}

\begin{remark}[What relabelling models]
\label{rem:relabel}
Two descriptions of the same agent that agree on all separations and
their costs are isomorphic. Renaming the agent's internal states,
re-indexing its memories, presenting its structure in a different
coordinate system---all are relabellings. An invariant is precisely a
feature of the agent that no such re-description can alter. Identity, if
it is anything, must be an invariant: it is what the agent is regardless
of how it is described.
\end{remark}

\subsection{The character invariant}

\begin{definition}[Internal residual; character invariant]
\label{def:char}
For the self-graph $\Self$, and any partition
$\mathcal{Q}=\{U_1,\dots,U_k\}$ of $\V$ into $k\ge2$ nonempty parts, the
\emph{internal residual} of the partition is the total inter-part
boundary cost
\[
\rho(\mathcal{Q})=\sum_{i<j}\wt(\cut(U_i,U_j)).
\]
The \emph{character invariant} of $\Agent$ is the minimum internal
residual over all such partitions,
\[
\Char(\Agent)=\min_{\mathcal{Q}:\,k\ge2}\rho(\mathcal{Q}),
\]
attained by finiteness of $\V$.
\end{definition}

The character invariant is the least total cost of splitting the agent
into pieces at all---the cheapest way to break the agent apart, measured
in the currency of its own separations. It is a property of the whole
internal structure, not of any one part.

\begin{theorem}[Identity is a conserved region]
\label{thm:identity}
Under \Cref{ax:finite,ax:present,ax:costly}:
\begin{enumerate}[label=\textup{(\roman*)},leftmargin=2.2em]
\item \textbf{Positive.} $\Char(\Agent)\ge\floorc>0$.
\item \textbf{Conserved.} $\Char$ is a weighted-graph invariant: if
$\phi:\Self\to\Self'$ is a weighted isomorphism then
$\Char(\Agent)=\Char(\Agent')$. In particular $\Char$ is fixed by every
relabelling of the agent's parts.
\item \textbf{Non-local.} $\Char$ is in general realised by a partition
into multi-part blocks, not by isolating a single part; it is a region of
the agent, never a point.
\end{enumerate}
\end{theorem}

\begin{proof}
(i) A partition with $k\ge2$ has, by connectedness (\Cref{ax:present}), at
least one crossing edge, of weight $\ge\floorc$ (\Cref{ax:costly}), so
$\rho(\mathcal{Q})\ge\floorc$; the minimum inherits the bound.

(ii) For a partition $\mathcal{Q}$ of $\V$, its image
$\phi(\mathcal{Q})=\{\phi(U_1),\dots,\phi(U_k)\}$ is a partition of $\V'$,
and since $\phi$ preserves edges and weights,
$\wt'(\cut'(\phi(U_i),\phi(U_j)))=\wt(\cut(U_i,U_j))$ for every pair;
summing, $\rho(\phi(\mathcal{Q}))=\rho(\mathcal{Q})$. As
$\mathcal{Q}\mapsto\phi(\mathcal{Q})$ is a bijection between the
nontrivial partitions of $\V$ and of $\V'$ (inverse $\phi^{-1}$), the
minima agree: $\Char(\Agent')=\Char(\Agent)$.

(iii) It suffices to exhibit an agent whose minimising partition is not a
singleton split. Take two triangles (each of edge cost $2$, so internally
expensive to cut) joined by a single edge of cost $\floorc=2$. Cutting
the joining edge separates the two triangles at total cost $2$; isolating
any single vertex costs at least $2+2=4$ (two triangle edges). So the
minimum residual is realised by the two-block partition, not by any
singleton. Hence in general the invariant is non-local.
\end{proof}

\begin{remark}[Interpretation: character, on which nothing depends]
\label{rem:character}
\Cref{thm:identity} licenses reading $\Char(\Agent)$ as the agent's
\emph{standing identity}: a feature conserved under all re-description,
positive (the agent is never nothing), and spread across the agent rather
than located in any single state. If one reads the agent's momentary
states as its surface conduct (\Cref{sec:behaviour}), the character
invariant is the bottom from which that conduct derives---what stays the
same about the character across all the different things it does. Two
agents may share a character invariant without being the same agent (two
isomorphic self-graphs in different worlds); the invariant is the
character \emph{type}, and a particular agent is a token of it. We use
only the graph statement.
\end{remark}

\begin{corollary}[Character is not a state]
\label{cor:not-a-state}
No single part of the agent carries its character: for the two-triangle
agent of the proof, no vertex $v$ has $b(\{v\})=\Char(\Agent)$. An agent's
identity therefore cannot be read off any one of its momentary states; it
is the standing relation among all of them.
\end{corollary}

\begin{proof}
Immediate from \Cref{thm:identity}(iii): the minimiser is a two-block
partition and every singleton has strictly larger boundary cost.
\end{proof}

% =====================================================================
\section{Behaviour: a gated walk, and the need for a drive}
\label{sec:behaviour}
% =====================================================================

Identity is standing; behaviour moves. We model the agent's conduct as a
walk on its self-graph---a sequence of parts it passes through---and show
that the graph alone does not determine the walk. Something more is
needed to say what the agent does next, and naming that something is the
route to purpose.

\begin{definition}[Behaviour as a walk; state]
\label{def:walk}
A \emph{behaviour} of length $\ell$ is a walk $W=(v_0,v_1,\dots,v_\ell)$
with $v_iv_{i+1}\in\E$ for each $i$. Its \emph{committed count} is
$\ell$, the number of acts taken. The \emph{state} at step $\ell$ is the
pair $\state_\ell=(v_\ell,\ell)$: the current part together with the
count of acts taken to reach it.
\end{definition}

\begin{lemma}[Behaviour is underdetermined by the graph]
\label{lem:branch}
If some part $v$ has degree $\ge2$, then from $v$ there are at least two
distinct one-step behaviours, and the self-graph alone does not determine
which is taken.
\end{lemma}

\begin{proof}
Degree $\ge2$ gives distinct neighbours $u_1,u_2$ and distinct admissible
steps $(v,u_1),(v,u_2)$. The triple $(\V,\E,\wt)$ carries no ordering of
the edges incident to $v$; it privileges neither. Hence the next step is
not a function of the self-graph.
\end{proof}

\begin{remark}[Branching is the generic case]
Any connected agent with $\abs{\V}\ge3$ that is not a bare path has a part
of degree $\ge2$; for an agent rich enough to hold many distinctions,
branching is normal. So underdetermination is not an edge case: the agent
almost always faces a choice its own structure does not settle.
\end{remark}

\begin{definition}[Drive field]
\label{def:drive}
A \emph{drive field} on $\Self$ is a map $\drive:\V\times\E\to[0,1]$
assigning to each part $v$ and incident separation $e\ni v$ a forward
weight $\drive(v,e)$ with $\sum_{e\ni v}\drive(v,e)\le1$ for each $v$. A
behaviour is \emph{$\drive$-admissible} if every step $v_iv_{i+1}$ has
$\drive(v_i,v_iv_{i+1})>0$. The drive \emph{selects} at $v$ when it gives
positive weight to exactly one incident separation.
\end{definition}

\begin{theorem}[Selection requires a drive]
\label{thm:drive}
Suppose some part of $\Agent$ has degree $\ge2$. Then a behaviour
determined step by step requires a drive field: no rule depending on
$\Self$ alone selects a unique next act at every branching part.
Conversely, any rule that does select a unique next act at each part
\emph{is} a selecting drive field.
\end{theorem}

\begin{proof}
\emph{Necessity.} At a branching part, \Cref{lem:branch} exhibits $\ge2$
admissible steps not distinguished by $\Self$. A rule that nonetheless
outputs one of them uses a datum beyond $(\V,\E,\wt)$; encode that datum
as $\drive(v,\cdot)$, positive on the chosen edge and zero elsewhere. This
is a drive field by \Cref{def:drive}. \emph{Converse.} If $\drive$ gives,
at each part, positive weight to exactly one incident edge, then from any
start the next admissible part is unique at every step, so $\drive$
determines a unique behaviour.
\end{proof}

\begin{remark}[Interpretation: motivation, on which nothing depends]
\label{rem:motivation}
The drive is the situation-dependent bias that decides what the agent
does next when its structure leaves the choice open. Read as
\emph{motivation}, it is what makes the same agent, in the same internal
state, act differently under different conditions: different drives are
different orderings over the same reachable interior. We use only the
graph statement (\Cref{thm:drive}); the reading is orientation.
\end{remark}

\begin{proposition}[Histories do not return]
\label{prop:nonreturn}
Along any behaviour the committed count is strictly increasing, so for
$i<j$ the states differ, $\state_i\neq\state_j$, even when $v_i=v_j$. A
revisited part is never a revisited state.
\end{proposition}

\begin{proof}
Each act increments the count by one, so counts are distinct for distinct
steps; since the state includes the count, $\state_i=(v_i,i)\neq
(v_j,j)=\state_j$ whenever $i\neq j$, regardless of the parts.
\end{proof}

\Cref{prop:nonreturn} is used in \Cref{sec:unreachable}: it is why an
agent that ``comes back to'' an earlier condition is not the earlier
agent, and why an authored copy is a new individual.

% =====================================================================
\section{Purpose: the fixed point of the drive}
\label{sec:purpose}
% =====================================================================

We now define purpose and prove it exists, is unique up to a natural
equivalence, and attracts. To speak of ``where behaviour settles'' we put
a mild, standard structure on the drive: we take it to be generated by a
potential, so that behaviour flows downhill and purpose is the bottom.
This is the one place we use continuous dynamics, and we keep it minimal.

\subsection{Behaviour as descent on a potential}

\begin{definition}[Behaviour space; drive potential]
\label{def:behspace}
Let the \emph{behaviour space} $\Beh$ be a compact convex subset of $\R^n$
coordinatising the agent's dispositions (for instance, for each part a
tendency-to-occupy in $[0,1]$, so $\Beh=[0,1]^{\abs\V}$). A \emph{drive
potential} is a continuously differentiable $\Phi:\Beh\to\R$ that is
$\lambda$-strongly convex for some $\lambda>0$: for all $x,y\in\Beh$,
\[
\Phi(y)\ \ge\ \Phi(x)+\nabla\Phi(x)^\top(y-x)+\tfrac{\lambda}{2}\norm{y-x}^2.
\]
The induced \emph{behaviour flow} is the gradient descent
$\dot x=-\nabla\Phi(x)$, with steps of the gated walk of
\Cref{sec:behaviour} realised as its discretisation.
\end{definition}

\begin{remark}[Why strong convexity is the honest minimal assumption]
\label{rem:convex}
We assume the drive has a single basin (one strongly convex well). This
is the assumption that an agent has \emph{a} purpose---one thing it is
drawn toward---rather than several competing ones. It is a modelling
choice, not a theorem: a multi-basin drive would give an agent several
candidate purposes and a selection problem among them. We make the
single-basin assumption explicit here and treat the multi-basin case as
out of scope (\Cref{sec:discussion}). Under it, existence, uniqueness, and
attraction are clean.
\end{remark}

\subsection{Purpose defined; existence and uniqueness}

\begin{definition}[Purpose]
\label{def:purpose}
A \emph{purpose} of $\Agent$ under drive potential $\Phi$ is a fixed point
of the behaviour flow: a point $x^\star\in\Beh$ with
$\nabla\Phi(x^\star)=0$ (an interior minimiser) or, at the boundary, a
Karush--Kuhn--Tucker point of $\min_{x\in\Beh}\Phi(x)$. The \emph{purpose
region} $\Purp$ is the set of such fixed points.
\end{definition}

\begin{theorem}[Purpose exists and is unique]
\label{thm:purpexist}
For a drive potential $\Phi$ on a compact convex $\Beh$, the purpose
region $\Purp$ is a single point $x^\star$: the flow has exactly one fixed
point, the unique minimiser of $\Phi$ on $\Beh$.
\end{theorem}

\begin{proof}
$\Phi$ is continuous on the compact set $\Beh$, so it attains a minimum;
strong convexity makes the minimiser unique (a strongly convex function on
a convex set has at most one minimiser). At an interior minimiser
$\nabla\Phi(x^\star)=0$; at a boundary minimiser the KKT conditions hold.
Conversely any fixed point of $\dot x=-\nabla\Phi$ in the interior has
$\nabla\Phi=0$ and, by convexity, is the global minimiser; a boundary
fixed point is a KKT point, again the minimiser by convexity. Hence
$\Purp=\{x^\star\}$.
\end{proof}

\begin{remark}[Uniqueness ``up to equivalence'']
\label{rem:equiv}
The minimiser $x^\star$ is unique as a point of $\Beh$. But behaviour is
observed only through the action it produces, and many nearby dispositions
produce the same action. If we quotient $\Beh$ by the equivalence ``yields
the same action,'' purpose is a \emph{cell}---a positive-diameter region
of dispositions all realising the purpose---and is unique as such. This
matches the fact (\Cref{thm:identity}(iii), \Cref{cor:not-a-state}) that
the agent's meaningful quantities are regions, not points: purpose is
unique, and it is a region.
\end{remark}

\subsection{Purpose attracts}

\begin{theorem}[Purpose is an attractor]
\label{thm:attractor}
Under a drive potential $\Phi$ with strong-convexity constant $\lambda$,
the purpose $x^\star$ is globally exponentially attracting on $\Beh$:
along the behaviour flow,
\[
\norm{x(t)-x^\star}\ \le\ e^{-\lambda t}\,\norm{x(0)-x^\star}.
\]
Behaviour perturbed away from purpose returns to it.
\end{theorem}

\begin{proof}
Let $V(x)=\tfrac12\norm{x-x^\star}^2$. Along $\dot x=-\nabla\Phi(x)$,
\[
\dot V=(x-x^\star)^\top\dot x=-(x-x^\star)^\top\nabla\Phi(x).
\]
Since $\nabla\Phi(x^\star)=0$ (or, at the boundary, the KKT inequality
holds in the descent direction), strong convexity gives the monotonicity
$(x-x^\star)^\top(\nabla\Phi(x)-\nabla\Phi(x^\star))\ge\lambda\norm{x-
x^\star}^2$, hence $\dot V\le-2\lambda V$. Grönwall's inequality yields
$V(t)\le e^{-2\lambda t}V(0)$, i.e. the stated bound.
\end{proof}

\begin{corollary}[Robust character]
\label{cor:robust}
A character defined by an identity $\Char(\Agent)$ and a purpose
$x^\star$ is robust: a disturbance that pushes conduct off purpose decays
at rate $\lambda$, and the agent resumes pursuing its purpose without
external correction. The larger $\lambda$, the more single-minded the
character; small $\lambda$ gives a character easily deflected and slow to
return.
\end{corollary}

\begin{proof}
Immediate from \Cref{thm:attractor}: the return rate is exactly the
strong-convexity constant $\lambda$ of the drive.
\end{proof}

\begin{remark}[Interpretation: goal, on which nothing depends]
The purpose $x^\star$ is the region of conduct the agent is drawn to and
returns toward---its \emph{goal} in the operational sense of the attractor
of its own motivation, not a symbolic target handed to it from outside.
An agent has a goal because its drive has a bottom, and it keeps the goal
under perturbation because the bottom attracts. We use only the dynamical
statements (\Cref{thm:purpexist,thm:attractor}).
\end{remark}

% =====================================================================
\section{Agents in society}
\label{sec:society}
% =====================================================================

An agent alone has identity and purpose. We now ask what happens when
agents act together. We show that a group acting as one has a single
shared drive (and that two drives mean two groups), that coordinated
purpose needs no shared identity, and that a crowd sharpens purpose. The
key device is that a society of agents is \emph{again} a graph, so the
whole apparatus applies one level up.

\subsection{Society as a quotient}

\begin{definition}[Society; quotient self-graph]
\label{def:society}
Let $\Agent_1,\dots,\Agent_m$ be agents, each with self-graph $\Self_i$,
embedded as disjoint blocks $U_1,\dots,U_m$ of a common vertex set
$\V^\ast=\bigsqcup_i\V_i$, with \emph{inter-agent separations} (edges
between blocks) recording the costs at which agents are held apart. The
\emph{society graph} $\Soc$ has one vertex per agent and an edge
$U_iU_j$, of weight $\wt(\cut(U_i,U_j))$, whenever some inter-agent
separation joins $\Agent_i$ to $\Agent_j$.
\end{definition}

\begin{lemma}[The society is an agent-shaped object]
\label{lem:society-graph}
If each inter-agent separation has cost $\ge\floorc$ and the society is
connected, then $\Soc$ is a finite connected weighted graph with floor
$\ge\floorc$: it satisfies \Cref{ax:finite,ax:present,ax:costly}.
\end{lemma}

\begin{proof}
$\Soc$ has $m<\infty$ vertices. Each edge weight is a sum of inter-agent
separation costs, each $\ge\floorc$, so $\ge\floorc>0$. Connectedness is
assumed. Hence the three axioms hold for $\Soc$.
\end{proof}

\begin{corollary}[Society has an identity]
\label{cor:society-identity}
By \Cref{thm:identity} applied to $\Soc$, a connected society carries its
own character invariant $\Char(\Soc)\ge\floorc>0$, conserved under
relabelling of \emph{which agent is which} and realised by no single
agent. A group has an identity that is not the identity of any of its
members.
\end{corollary}

\begin{proof}
\Cref{lem:society-graph} makes $\Soc$ satisfy the axioms;
\Cref{thm:identity} then gives the invariant, its conservation, and its
non-locality, now read at the level of agents.
\end{proof}

\subsection{One group, one drive}

\begin{definition}[Coherent society]
\label{def:coherent}
A society $\Soc$ is \emph{coherent} if it carries a single selecting drive
field (\Cref{def:drive}) on $\Soc$ that picks, at each branching agent, a
unique next agent to act.
\end{definition}

\begin{theorem}[One group, one drive]
\label{thm:onedrive}
Let $\Soc$ have an agent of degree $\ge2$. Then:
\begin{enumerate}[label=\textup{(\roman*)},leftmargin=2.2em]
\item selecting which agent acts next requires a society-level drive
(\Cref{thm:drive} applied to $\Soc$); and
\item if two selecting drives both render $\Soc$ coherent, they agree
wherever either selects. A society admitting two genuinely different
selecting drives is not one coherent society but two, partitioned into the
sub-societies on which each drive is the selector.
\end{enumerate}
\end{theorem}

\begin{proof}
(i) is \Cref{thm:drive} for the graph $\Soc$ (\Cref{lem:society-graph}).
(ii) At each branching agent $U$, coherence means a unique next agent is
determined; a selecting drive puts positive weight on exactly the edge to
that agent and zero elsewhere, so any two selecting drives have the same
support at $U$ and select the same successor. Thus they agree wherever
either selects. If instead two drives select different successors at some
agent, group ``shares a selected successor'' into blocks within which a
single drive selects consistently; these blocks are distinct
sub-societies, so the original was a union of coherent societies, not one.
\end{proof}

\begin{remark}[Interpretation: shared purpose and schism]
\label{rem:schism}
Read the society-level drive as the group's \emph{shared purpose}: the
field that selects which member acts, toward what. \Cref{thm:onedrive}
then says a coherent group has exactly one such field---one purpose---and
that two purposes are the signature of a group that has really split. The
society's own attractor (\Cref{thm:attractor} applied to a drive potential
on the society's behaviour space) is the group's goal, standing to the
group as an individual's purpose stands to the individual.
\end{remark}

\subsection{Coordination without shared identity}

The strongest social claim is that agents need share \emph{nothing}
internal to pursue a common purpose. We make ``common purpose'' live in
the space of outcomes, not in any agent's head.

\begin{definition}[Outcome space; outcome purpose]
\label{def:outcome}
Let $(\Beh^\ast,\dist)$ be a common \emph{outcome space}: a metric space
of observable results, shared by all agents, with each agent equipped with
a map $A_i$ from its own behaviour space $\Beh_i$ to $\Beh^\ast$ (the
outcome its conduct produces). A \emph{shared purpose} is a region
$\Purp^\ast\subseteq\Beh^\ast$ with positive tolerance
$\tau(\Purp^\ast)>0$ (there are outcomes within $\tau$ of it). Agent
$\Agent_i$ \emph{attains} $\Purp^\ast$ if its purpose $x^\star_i$ maps
near it: $\dist(A_i(x^\star_i),\Purp^\ast)\le\tau(\Purp^\ast)$.
\end{definition}

\begin{theorem}[Coordination without shared internals]
\label{thm:coord}
Agents with wholly disjoint self-graphs, disjoint behaviour spaces, and
different drives can all attain the same shared purpose
$\Purp^\ast$. Formally, if for each $i$ the agent's purpose satisfies
$\dist(A_i(x^\star_i),\Purp^\ast)\le\tau(\Purp^\ast)$, then every agent
attains $\Purp^\ast$, regardless of whether the $\Self_i$, $\Beh_i$, or
drives share any content.
\end{theorem}

\begin{proof}
The condition to attain $\Purp^\ast$ is a statement purely about the image
$A_i(x^\star_i)$ in the common outcome space $\Beh^\ast$: it holds for
each $i$ by hypothesis. Nothing in the condition refers to the internal
structure $\Self_i$, the behaviour space $\Beh_i$, or the drive of any
agent, nor to any relation between those of different agents. Hence the
disjointness of internals is compatible with joint attainment: all agents
land their outcomes within tolerance of the same region.
\end{proof}

\begin{remark}[Why this matters for characters]
\Cref{thm:coord} says a crowd of characters can pursue a shared purpose
without any shared mind, vocabulary, or belief: three onlookers keep a
safe distance from a dangerous animal though one is an expert, one a
novice, one a child, because the shared purpose (``stay clear'') lives in
outcome space and each reaches it by its own internal route. For synthetic
agents this means one need not give a party of characters a common
world-model to make them act coherently; overlapping outcomes suffice.
\end{remark}

\subsection{Crowds sharpen purpose}

\begin{definition}[Independent attainment]
\label{def:indep}
Agents $\Agent_1,\dots,\Agent_n$ attain a shared purpose $\Purp^\ast$
\emph{independently} if the events ``$\Agent_i$ fails to bring the
collective outcome within tolerance'' are independent, with per-agent
failure probability $q_i\in[0,1]$.
\end{definition}

\begin{theorem}[Crowd sharpening]
\label{thm:crowd}
Under independent attainment with a parallel (at-least-one-succeeds)
reading of the collective outcome, the collective failure probability is
$\prod_{i=1}^n q_i$, which is non-increasing in $n$ and strictly
decreasing whenever an added agent has $q_i<1$. A crowd attains the shared
purpose at least as sharply as, and generically more sharply than, any
sub-crowd; and if infinitely many agents keep $q_i$ bounded below $1$, the
collective failure tends to $0$.
\end{theorem}

\begin{proof}
By independence, the collective fails iff every agent fails, with
probability $\prod_i q_i$. Adding an agent multiplies by $q_{n+1}\le1$, so
the product is non-increasing, and strictly decreasing if $q_{n+1}<1$. If
$q_i\le q<1$ for infinitely many $i$, then $\prod_i q_i\le q^{\,k}\to0$ as
the number $k$ of such agents grows. (This is the Borel--Cantelli-type
dichotomy: $\prod_i q_i\to0$ iff $\sum_i(1-q_i)$ diverges, e.g. when the
$q_i$ stay bounded below $1$.)
\end{proof}

\begin{remark}[Emergent group competence]
\Cref{thm:crowd} gives, for free, that a group of individually mediocre
characters can be collectively reliable at a shared purpose. The
sharpening is multiplicative in the number of participants; a
few strongly-attaining agents or many weakly-attaining ones both drive the
collective failure down. This is the quantitative form of ``the crowd gets
it right even when no one is sure.''
\end{remark}

% =====================================================================
\section{Identity is unreachable from outside; copies are distinct}
\label{sec:unreachable}
% =====================================================================

Finally we show two things that give characters their felt autonomy: an
outside observer cannot resolve an agent's identity to a point, and an
authored copy of a character is always a distinct individual.

\begin{definition}[Interrogation]
\label{def:interrogation}
An \emph{interrogation} of $\Agent$ is a finite sequence of separations an
outside party adds to $\Agent$'s boundary---finitely many distinctions
drawn between $\Agent$ and the rest of some world---each tightening how
$\Agent$ is told apart from things it is not.
\end{definition}

\begin{theorem}[Identity is unreachable to a point]
\label{thm:unreachable}
No finite interrogation resolves an agent's identity to a point. Each
added separation tightens the agent's \emph{exterior}---what it is
distinguished from---but leaves its character invariant $\Char(\Agent)$
unchanged, since $\Char$ is an invariant of the agent's own structure and
adding boundary edges is a change of the exterior, not of the internal
partition minimising $\rho$. Hence after any finite interrogation the
identity remains a region of size $\ge\floorc$, never a point.
\end{theorem}

\begin{proof}
$\Char(\Agent)$ is defined (\Cref{def:char}) by minimising internal
residual over partitions of $\V$, the agent's own parts. An interrogation
adds edges between $\V$ and vertices outside $\Agent$; it does not alter
$\Self=(\V,\E,\wt)$ and so does not change any $\rho(\mathcal Q)$ for
partitions $\mathcal Q$ of $\V$. Thus $\Char(\Agent)$ is invariant under
interrogation, and by \Cref{thm:identity}(i),(iii) it remains
$\ge\floorc>0$ and non-local---a region. What the interrogation reduces is
the exterior ambiguity (which things $\Agent$ might be confused with), not
the interior invariant. Since no finite interrogation drives the interior
to a point, identity is unreachable to a point from outside.
\end{proof}

\begin{corollary}[To know a character exactly is impossible from outside]
An outside party can make an agent's identity arbitrarily well-separated
from other things without ever possessing the identity itself; the
interior invariant is approached from outside and never delivered. This is
the precise sense in which a well-made character retains an interior the
player never fully reaches.
\end{corollary}

\begin{proof}
Immediate from \Cref{thm:unreachable}: interrogation moves only the
exterior; the interior invariant is conserved.
\end{proof}

\begin{theorem}[An authored copy is a distinct individual]
\label{thm:copy}
Let $\Agent$ have behaviour history of committed count $m>0$, and let
$\Agent'$ be authored with self-graph isomorphic to $\Agent$'s but started
fresh at count $0$. Then $\Agent'\neq\Agent$: their states differ in the
committed count, and $\Agent'$ satisfies the predicate ``authored with
empty history'' that $\Agent$ does not. Resemblance of self-graph does not
make the copy identical to its original.
\end{theorem}

\begin{proof}
A state includes the committed count (\Cref{def:walk}); $\Agent$ is at
count $m>0$, $\Agent'$ at count $0$, so their states differ. By
\Cref{prop:nonreturn} no behaviour lowers the count, so $\Agent'$ cannot
acquire $\Agent$'s history without taking the intervening acts, which
would make it a further-distinct individual. And $\Agent'$ is individuated
by a distinction $\Agent$ lacks (empty prior history). Hence
$\Agent'\neq\Agent$.
\end{proof}

\begin{remark}[Authored, yet its own]
\Cref{thm:unreachable} and \Cref{thm:copy} together say an authored
character is genuinely its own: its interior is not exhaustible by outside
interrogation, and duplicating its structure yields a new individual with
its own history, not the same one again. One authors the standing
structure and the drive; one does not author the lived history, which the
monotone count individuates.
\end{remark}

% =====================================================================
\section{Numerical study}
\label{sec:validation}
% =====================================================================

The results are proved from the axioms and need no computation; we
nonetheless exhibit them on explicitly constructed agents, to display
their quantitative shape and to check the derivations by direct
enumeration.

\subsection{Method}

An agent is realised as a finite connected weighted graph with all edge
weights drawn at or above a floor $\floorc>0$. On each agent we compute,
by exact enumeration over subsets and partitions: every boundary cost
$b(U)$; the realised floor $\floorc_\Self$; the character invariant
$\Char(\Agent)$ and a minimising partition; and, for a drive potential
$\Phi$ chosen strongly convex on $\Beh=[0,1]^{\abs\V}$, the purpose
$x^\star$ and the empirical return rate along the behaviour flow from
random initial conditions. A \emph{relabelling} is a random
weight-preserving permutation of parts; we recompute $\Char$ after each to
test invariance. A \emph{society} is built by joining several agents with
inter-agent edges $\ge\floorc$ and computing $\Char(\Soc)$ and its
minimising agent-partition. A \emph{crowd} of $n$ independent attainers
with per-agent failure $q_i$ is checked against the predicted collective
failure $\prod_i q_i$.

\subsection{Predictions checked}

\begin{enumerate}[label=\textbf{(V\arabic*)},leftmargin=*]
\item \textbf{Floor.} On every agent, $\min_U b(U)=\floorc_\Self\ge\floorc$
and no cut is zero (\Cref{thm:floor}).
\item \textbf{Identity is conserved.} $\Char(\Agent)$ is unchanged to
machine precision under every relabelling (\Cref{thm:identity}(ii)).
\item \textbf{Identity is non-local.} On the two-triangles family the
minimiser of $\Char$ is the two-block partition, never a singleton
(\Cref{thm:identity}(iii)).
\item \textbf{Purpose exists, is unique, attracts.} The behaviour flow has
a single fixed point, reached from every initial condition with empirical
rate matching the strong-convexity constant $\lambda$
(\Cref{thm:purpexist,thm:attractor}).
\item \textbf{One group, one drive.} A society built from two internally
coherent sub-groups joined by a single edge admits two selecting drives
exactly when the joining edge is cut, i.e. exactly when it is two
societies (\Cref{thm:onedrive}).
\item \textbf{Coordination without shared internals.} Agents with disjoint
self-graphs and different drives, but outcome maps landing near a common
$\Purp^\ast$, all attain it (\Cref{thm:coord}).
\item \textbf{Crowd sharpening.} Collective failure equals $\prod_i q_i$
and decreases geometrically in $n$ for $q_i\le q<1$ (\Cref{thm:crowd}).
\item \textbf{Copies are distinct.} A fresh copy differs from its original
in committed count at every step (\Cref{thm:copy}).
\end{enumerate}

We report that on a sweep of randomly generated agents (varying part
count, floor, edge density, and weight spread), together with the
two-triangle witness and a battery of societies and crowds, every
predicted inequality, invariance, and identity holds: no cut is sharp;
$\Char$ is invariant under every relabelling; its minimiser is non-local
on the witness family; the behaviour flow has a unique attractor with the
predicted return rate; societies split into two drives exactly at the
severing edge; disjoint agents co-attain a common outcome purpose; and
collective failure follows $\prod_i q_i$. (The suite is deterministic
given its inputs; figures and a reproducible harness accompany the source.)

% =====================================================================
\section{Discussion}
\label{sec:discussion}
% =====================================================================

\subsection{What has and has not been established}

We have shown that any bounded agent satisfying finiteness, presence, and
costly individuation has a conserved, non-local, positive
\emph{identity}---its character invariant---and, once equipped with a
single-basin drive, a unique \emph{purpose} that attracts. In society, a
coherent group has one shared drive, coordinated purpose needs no shared
internals, and crowds sharpen purpose; and identity is unreachable from
outside while copies are distinct. These are theorems from the stated
axioms.

We have \emph{not} shown that a single-basin drive is the right model of
motivation for every character; the strong-convexity assumption
(\Cref{rem:convex}) is a modelling choice that buys clean existence and
attraction at the price of ruling out characters torn between competing
purposes. Nor have we claimed that human or animal identity is a finite
graph; the axioms are meant to hold of \emph{synthetic} bounded agents,
where they are design commitments one can honour by construction.

\subsection{Open directions}

The natural extensions are: (i) multi-basin drives, giving characters with
competing purposes and a selection dynamics among them; (ii) drives that
change over time as the agent's history accumulates, coupling purpose back
to the monotone count of \Cref{prop:nonreturn}; and (iii) the mechanics of
how an agent processes an incoming interaction and decides whether it is
worth responding to---which is the subject of the companion paper, where
the identity and purpose defined here become the standard against which an
interaction's relevance is measured.

\subsection{Closing}

Identity is what an agent conserves; purpose is where it settles. A
character is the pair. The definitions are region-valued, not
point-valued, which is why a good character has an interior an observer
never exhausts and a purpose it keeps under disturbance; and the same
construction, applied one level up, gives a society its own identity and
its own single purpose. Nothing here needs to be sold: it is a set of
theorems about finite agents, and the readings of them as characters are
offered only as orientation.

\vspace{0.6em}
\hrule
\vspace{0.4em}
\noindent\emph{Identity is the conserved interior; purpose is the fixed
point of the drive; a society is an agent one level up.}

\bibliographystyle{plain}
\bibliography{references}

\end{document}
